\chapter{Doctoral Entrance Problems}
\label{ch:problems}

The problems below are reproduced in French and routed by individual mathematical content.
\newpage

\section{Université Badji Mokhtar --- Annaba}
\begin{problem}{(Annaba) — Épreuve N°02}{ds:annaba-2015-e2-p7}
Considérons le système différentiel réel suivant
$$\begin{cases}x'=-ax+ay,\\ y'=bx-y-xz,\\ z'=-cz+xy,\end{cases}$$
où $c>0$, $a>c+1$ et $b>0$, appelé système différentiel de Lorenz.
\begin{enumerate}\item Vérifier que dans le cas $b<1$ le seul point d'équilibre est la solution nulle.\item Déterminer les points d'équilibre de ce système différentiel dans le cas $b>1$.\item Étudier la stabilité de la solution nulle.\end{enumerate}
Soit dans $\mathbb{R}^3$ la courbe d'équations $x=\cos t$, $y=\sin t$, $z=bt$ où $b$ est une constante.
\begin{enumerate}\setcounter{enumi}{3}\item Faire un dessin.\item Déterminer le trièdre de Frenet-Serret $(T,N,B)$.\item Démontrer que la courbure et la torsion sont constantes en tout point.\end{enumerate}
\source{\emph{\small Université Badji Mokhtar - Annaba, 2015, Épreuve~n\textsuperscript{o}~02.}}\end{problem}
\newpage

\section{Université Frères Mentouri --- Constantine 1}
\begin{problem}{(Constantine 1) — Épreuve N°11}{ds:constantine1-2015-e11-p1}
Supposons que $f$ et $g$ sont deux fonctions réelles continues dans l'intervalle $a\le t\le b$, et que $f'$ existe sur cet intervalle, telles que
$$f'(t)\le f(t)g(t),\qquad a\le t\le b.\tag{1}$$
Posons alors $q(t)=\exp\left(-\int_a^tg(s)\,ds\right)$.
\begin{enumerate}\item Calculer la dérivée $q'(t)$.\item Montrer que la fonction $F(t)=f(t)q(t)$ est non croissante.\item Montrer que pour tout $t\in[a,b]$, on a
$$f(t)\le f(a)\exp\left(\int_a^tg(s)\,ds\right).$$\end{enumerate}
\source{\emph{\small Université Frères Mentouri - Constantine 1, 2015, Épreuve~n\textsuperscript{o}~11.}}\end{problem}
\newpage
\begin{problem}{(Constantine 1) — Épreuve N°01}{ds:constantine1-2013-e1-p4a}
Résoudre par la transformation de Laplace $\mathcal L$ le problème $(P)$ :
$$(P):\begin{cases}y''+4y'+3y=\delta(t),\\ y(0)=0,\quad y'(0)=0,\end{cases}$$
où $\delta$ est la mesure de Dirac et $\mathcal L[f](s)=\int_0^{+\infty}e^{-st}f(t)\,dt$.
\source{\emph{\small Université Frères Mentouri - Constantine 1, 2013, Épreuve~n\textsuperscript{o}~01.}}\end{problem}
\newpage
\begin{problem}{(Constantine 1) — Épreuve N°01}{ds:constantine1-2013-e1-p4b}
Soit le problème à valeurs initiales suivant :
$$(P):\begin{cases}y''+4y'+5y=\delta(t),\\ y(0)=0,\quad y'(0)=0,\end{cases}$$
où $\delta(t)$ est la fonction de Dirac. Résoudre le problème $(P)$ en utilisant la transformation de Laplace $\mathcal L$. La transformation de Laplace $\mathcal L$ de $f(t)$, notée $F(s)=[\mathcal Lf](s)$, est définie par $F(s)=[\mathcal Lf](s)=\int_0^\infty e^{-st}f(t)\,dt$.
\source{\emph{\small Université Frères Mentouri - Constantine 1, 2013, Épreuve~n\textsuperscript{o}~01.}}\end{problem}

\newpage
\section{Université Djilali Liabès --- Sidi Bel Abbès}
\begin{problem}{(Sidi Bel Abbès) — Épreuve N°01}{ds:sba-2021-e1-p2}
Soit $A : ]0, +\infty[ \to M_n(\mathbb{R})$ une application continue. On considère l'équation différentielle linéaire
$$
x'(t) = A x(t), \qquad t \in ]0, +\infty[, \tag{3}
$$
et on note $R(t, t_0)$ sa résolvante.
\begin{enumerate}
\item Montrer que $S(t, t_0) = R(t_0, t)^T$ est la résolvante de l'équation différentielle
$$z'(t) = -A(t)^T z(t).$$
\item Dans cette équation on choisit
$$
A(t) = \begin{pmatrix} 2t + \frac{1}{t} & 0 & \frac{1}{t} - t \\ t - \frac{1}{t} & 3t & t - \frac{1}{t} \\ \frac{2}{t} - 2t & 0 & \frac{2}{t} + t \end{pmatrix}.
$$
Montrer que $A(t)$ possède une base de vecteurs propres indépendante de $t$. En déduire la résolvante $R(t, t_0)$ de l'équation (3).
\item Résoudre à l'aide des questions précédentes le système
$$
\begin{cases} x_1' = -(2t + \frac{1}{t}) x_1 + (\frac{1}{t} - t) x_2 + 2(t - \frac{1}{t}) x_3, \\ x_2' = -3t\, x_2, \\ x_3' = (t - \frac{1}{t}) x_1 + (\frac{1}{t} - t) x_2 - (\frac{2}{t} + t) x_3. \end{cases}
$$
\end{enumerate}
\source{\emph{\small Université Djilali Liabès - Sidi Bel Abbès, 2021, Épreuve~n\textsuperscript{o}~01.}}\end{problem}
\newpage
\begin{problem}{(Sidi Bel Abbès) — Épreuve N°03}{ds:sba-2021-e3-p1}
Une population constante $N$ est infectée par une épidémie ; le modèle décrivant l'évolution des compartiments $S$ (sains, susceptibles) et $I$ (infectés) est
$$S'=-\frac{\beta}{N}SI+(d+\lambda)I,\quad t\in\mathbb{R}_+,\tag{1}$$
$$I'=\frac{\beta}{N}SI-(d+\lambda)I,\quad t\in\mathbb{R}_+,\tag{2}$$
où $d$ est le taux de mortalité, $\lambda$ le taux de guérison et $\beta$ le taux d'infection.
\begin{enumerate}
\item Quel est le type du modèle (1)-(2) ?
\item En déduire le taux de naissance de la population $N$.
\item Montrer que le système (1)-(2) admet une unique solution positive globale, pour toute condition initiale $N(0)\ge0$, $I(0)\ge0$.
\item Écrire l'équation différentielle (2) en fonction de $I$ et des paramètres du modèle.
\item Calculer $I$, et en déduire la limite de $I(t)$ quand $t\to+\infty$.
\end{enumerate}
\source{\emph{\small Université Djilali Liabès - Sidi Bel Abbès, 2021, Épreuve~n\textsuperscript{o}~03.}}\end{problem}
\newpage
\begin{problem}{(Sidi Bel Abbès) — Épreuve N°12347}{ds:sba-2026-e12347-p1}
Pour le système linéaire $\dot{x}=Ax$, on considère les matrices
$$
A_1=\begin{pmatrix}2&1\\0&-1\end{pmatrix},\quad A_2=\begin{pmatrix}-1&2\\-2&-1\end{pmatrix},\quad A_3=\begin{pmatrix}0&1\\-1&0\end{pmatrix},\quad A_4=\begin{pmatrix}-2&0\\0&-3\end{pmatrix}.
$$
\begin{enumerate}
\item Calculer les valeurs propres de chaque matrice.
\item Classifier chaque équilibre (nœud, foyer, col, centre).
\item Préciser la stabilité de chaque équilibre.
\item Esquisser le portrait de phase pour $A_2$ et $A_4$.
\end{enumerate}
\source{\emph{\small Université Djilali Liabès - Sidi Bel Abbès, 2026, Épreuve~n\textsuperscript{o}~12347.}}\end{problem}
\newpage
\begin{problem}{(Sidi Bel Abbès) — Épreuve N°12347}{ds:sba-2026-e12347-p2}
Bien que discret, ce modèle illustre les concepts de la dynamique hôte-parasitoïde. On étudie ici la version continue approchée :
$$
\dot H=rH-aHP,\qquad \dot P=caHP-dP,
$$
où $H$ désigne les hôtes, $P$ les parasitoïdes, $a$ le taux de parasitisme, $c$ l’efficacité de conversion, $r$ le taux de croissance des hôtes et $d$ le taux de mortalité des parasitoïdes, avec $a,c,r,d$ strictement positifs.
\begin{enumerate}
\item Calculer les équilibres et comparer avec Lotka--Volterra.
\item Analyser la nature de l’équilibre (centre, foyer, nœud ?).
\item Ajouter une densité-dépendance pour les hôtes :
$$
\dot H=rH\left(1-\frac{H}{K}\right)-aHP.
$$
Étudier la stabilité.
\item Interpréter biologiquement dans le contexte de la lutte biologique.
\end{enumerate}
\source{\emph{\small Université Djilali Liabès - Sidi Bel Abbès, 2026, Épreuve~n\textsuperscript{o}~12347.}}\end{problem}
\newpage
\begin{problem}{(Sidi Bel Abbès) — Épreuve N°12347}{ds:sba-2026-e12347-p3}
Le modèle SEIR avec délai discret représentant la période d’incubation est
$$
\dot S=\Lambda-\beta SI-\mu S,
$$
$$
\dot E=\beta SI-\beta S(t-\tau)I(t-\tau)e^{-\mu\tau}-\mu E,
$$
$$
\dot I=\beta S(t-\tau)I(t-\tau)e^{-\mu\tau}-(\gamma+\mu)I,
$$
$$
\dot R=\gamma I-\mu R.
$$
\begin{enumerate}
\item Calculer $R_0$ et l’équilibre sans maladie $E_0$.
\item Montrer que $R_0$ décroît avec $\tau$.
\item Calculer l’équilibre endémique $E^*=(S^*,E^*,I^*,R^*)$ pour $R_0>1$.
\item Analyser qualitativement l’effet du retard sur la dynamique transitoire.
\end{enumerate}
\source{\emph{\small Université Djilali Liabès - Sidi Bel Abbès, 2026, Épreuve~n\textsuperscript{o}~12347.}}\end{problem}
